\section{Theoretical Framework}\label{sec:theory}

\subsection{Theoretical Background}\label{theo_background}
    Six member states of the European Union have retained their own currencies instead of adopting the euro.
The theory of optimum currency area address whether adopting same currency carries a cost (\cite{mundell_theory_1961}).
Mundell posed the question of whether a group of countries should share the same currency and argued currency boundaries are
an economic question rather than political. 
Pointed out that the shared currency removes the transaction costs of conversion and eliminates the future exchange rate uncertainty.
The benefit of monetary both larger the intensity of trade between country and lower the cost of cross-border trade and investment.
While at the same time it showed the other side of adopting common currency, 
losing the two instruments where are the independency of monetary policy and the ability to adjust the nominal exchange rate.
When symmetric shock applied, members of common currency tend to react the same and the common policy response is broadly appropriate for all members since it impose a much smaller cost.
In the other hand, the asymmetric shocks impose a larger cost where the adjustment must occur through wages, prices and employments. 
It showed slower and costly effects because single policy can not suit countries needing opposite responses.

\vspace{\baselineskip}

Mundell (\citeyear{mundell_theory_1961}) proposed factor mobility as the criterion to judge whether the benefits outweigh the cost.
When labor can freely move between country migration from weaker demand to stronger demand would absorbs the asymmetric shock replace unemployment.
Hence, the exchange rate do not needed for adjustment.
When the cost of giving up the exchange rate is small in regions of high factor mobility and large in limited mobility.
Since there are no operating adjustment must occur through employment.
Mckinnon (\citeyear{mckinnon_optimum_1963}) argued the openness as an alternative criterion. 
When most of the consumption is imported in a very open economy, the prices are determined by partner's currency.
When the currency depreciates and rapid increase in import prices leading the domestic prices and wages rise in the same direction.
Therefore, the gain from the cheaper exports do not last.
The exchange rate become less useful to open economic than closed one. Giving up the their own currency cost less.
Another perspective showing diversification as third condition, arguing economy that's producing diversity of goods is less exposed to shocks, which reduce the need of the exchange rate to response (\cite{kenen_theory_1969}).

\vspace{\baselineskip}

Those conditions that mentioned above were intended to be checked before forming a currency union.
Frankel and Rose (\citeyear{frankel_endogeneity_1998}) argued that those are endogenous to the decision.
Adopting a common currency increase trade between countries which they tend to experience similar business cycle.
Therefore, the country failed the condition before might satisfy them after joining the currency union.
It showing the trades run in both directions which we should consider the reported result can be association rather than just a effect.
The condition we discussed only concern the decision of giving up a nation currency. 
And we will consider what a country can gain or loss by keeping their own currency.

\vspace{\baselineskip}

When demand of country's exports decrease, their goods become cheaper relative to trade partners. 
Friedman (\citeyear{milton_methodology_1953}) posed two scenarios can happen where the domestic wages and prices decrease, but the general wage levels resist.
The adjustment will be slow and happen through unemployment.
Alternatively, the currency can depreciate showing the same immediate change in relative prices.
Friedman argued it was easier to allow one price to move than adjust.   
With this argument the flexible exchange rate works as a main idea.

\vspace{\baselineskip}

We can take Friedman's argument as an implication for how exchange rate volatility should be interpreted.
When exchange rate acting as the adjustment function the volatility we can visibly observed in the function which can be explain just the adjustment work.
A currency moving according to a shock not necessarily a cost, it could be absorbing shocks.
We can consider volatility and macroeconomic instability are different where we cannot interpret the currency movement and economic performance without knowing what the currency is responding to.

\vspace{\baselineskip}

Against the argument we mentioned above that exchange rate movement can be treated as a cost.
Exchange rate volatility might discourage trade through uncertainty.
We consider a firm that agree a price in foreign currency and paid months later.
The domestic currency value of the payment is uncertain.
When the firm is not a risk taker will take certain payment over uncertain payment with same expected size.
The uncertain part of the payment make the transaction less attractive and trades that would be profitable will lost the gain they generated.
Even a firm with hedging strategy enter the forward contract with fixed exchange rate would not remove the problem.
There will be a cost in the contract and limited market for smaller currency contracts.
Consider a Polish's firm exporting to the euro area pays higher price for protection and gain less than a local firm operating in major currency.
The evidence that volatility actually reduce trade is weaker than the theoretical argument.
Clark et al. (\citeyear{clark_new_2004}) examined trade among 178 countries and find no robust evidence that exchange rate volatility has large negative effect on trade.
The volatility is not the major policy concern.
This does not mean the mechanism could not hold. 
It is difficult to measure uncertainty so the real effect may hard to detect.
This matter more to the countries that trade heavily with one partner instead size of the cost.
When a country's trade conducted with euro area larger share of the economy affected by the uncertainty of exchange rate.

\vspace{\baselineskip}

The two mechanisms we mentioned above point in opposite directions.
Exchange rate flexibility absorbs a shock supporting growth and at the same time exchange rate volatility increase the cost of trade.
Both mechanisms operate in floating currency where it can be observed as their net effect.
The theory suggests the balance between them depends on how exposed a country is to currency.
McKinnon's condition implies when there's more intensively an economy's trades with euro area the less value an independent exchange rate has.
Other argument implies an economy exposed to larger exchange rate uncertainty when they traded more with euro area.
Two arguments show the same conclusion where trade with euro area increase lead a flexible currency less useful and costly.
The theory predicts the effect of exchange rate volatility on growth is not constant.
But it decreases when country's trade with euro area increases.
It does not predict the effect of any level of trade integration and where two mechanisms offset another.
A further point from adjustment mechanism.
Flexible exchange rate absorbs shocks only when there are shocks can be absorb.
Therefore, in calm periods the relationship of mechanism is unlikely to be detectable in the data.
We will discuss both questions in the empirical part of the paper.


\subsection{Literature Review}\label{subsec:lit_review} 

The relationship between exchange rate volatility and economic growth has remained unsettled for decades with theory and evidence pointing in different directions (\cite{morina_effect_2020}).
Theory shows negative effect on volatility by increasing cost of international transaction.
However, the estimated effects have varied with countries and different periods with the measure of volatility (\cite{clark_new_2004}).
Our focus here is the effect of volatility on growth and trade with behavior of the non-euro EU country.

\vspace{\baselineskip}

The study by Morina et al. (\citeyear{morina_effect_2020}) examines fourteen Central and Eastern European Countries over 2002 to 2018.
In this study volatility is measured from real effective exchange rate as a standard deviation and z-score, enters a panel growth regression.
They estimate the model with country fixed effects having rejected random effects on basis of Hausman test.
The control variables are government expenditure, gross fixed capital, inflation, trade openness and domestic credit to private sector.
Their analysis conducts in two stages, estimates the direct effect of the volatility on growth and interacting volatility with real interest rate, inflation and trade openness.
They found a significant negative effect of volatility on growth and positive effect of trade openness.
The interaction of real interest rate and inflation are found to reduce the negative effect of volatility.
Interaction between volatility and trade openness shows negative and insignificant at the ten percentage.
This trade interaction which is most relevant to the study is the least conclusive in result.
The measurement of volatility matters for the results.
Morina et al. (\citeyear{morina_effect_2020}) measures their volatility as average across a whole period. 
The calm and volatile periods are treated the same. 
It computed against a group of trading partners instead of a single currency.

\vspace{\baselineskip}

Other studies showed mixed evidences. Janus and Riera-Crichton (\citeyear{thorsten_pdf_2015}) show that a growth decrease might destabilize the exchange rate when foreign capital exits the economy.
They address real effective exchange rate with volatility of commodity terms of trade.
Estimate the model on an annual panel of OECD countries from 1980 to 2011.
It shows negative volatility with growth and one standard deviation fall in volatility associated with around two percentage points of additional growth.
They reports an important finding here. The Great Recession increased exchange rate volatility in the non-euro countries by 38 to 189 percent of the sample mean.
They used Denmark and Sweden as non-euro sample. This suggests exchange rate differ inside and outside euro area mainly during financial stress periods.

\vspace{\baselineskip}

There are several studies focus on how the currency move against the euro.
They indicate a measurement convention where object of study is exposure to a certain currency.
It measures exchange rates against that currency instead of a trade-weighted average.
Orlowski (\citeyear{orlowski_co-movements_2016}) studies co-movement between non-euro EU currencies and euro from 2000 to 2015.
Orlowski finds high co-movement for Hungarian forint, Czech koruna, Polish zloty, Swedish krona and Danish krone while Romanian leu showing little co-movement.
The co-movement rose during periods of financial crisis with Czech, Polish and Hungarian currencies (\cite{chomicz-grabowska_financial_2020})
Kobor and Szekely (\citeyear{kobor_foreign_2004}) find volatility in these countries against the euro is subject to regime switching.
It shows cross-correlations higher when both currencies are in the high volatility regime.

\vspace{\baselineskip}

Studies of volatility and growth tend to measure volatility by summary statistic computed by annual data.
While condition volatility models have been applied in studies of exchange rate movement rather than macroeconomic factors.
Bollerslev (\citeyear{bollerslev_generalized_1986}) generalized Engle (\citeyear{engle_autoregressive_1982}) study and introduced GARCH model which conditional variance of a series depends on recent financial shocks and variance.
This model was developed to response the features of financial data that fixed window standard deviation cannot represent.
For example periods of high volatility can follow even higher volatility.
So this model can estimate volatility for every observation in series instead of single period.
Latief and Lefen (\citeyear{latief_effect_2018}) estimate GARCH (1,1) and threshold GARCH (1,1) models with an AR(1) mean equation on daily exchange rate returns for seven South Asian countries over 1995 to 2016.
They use the result volatility series in panel regressions on international trade and foreign direct investment.
The volatility is positive with trade in Bhutan, the Maldives and Nepal and negative in Pakistan.

\vspace{\baselineskip}

Both theoretical and empirical literatures could not establish a clear relationship between exchange rate volatility and economic growth.
Theory shows two mechanisms working in opposite directions.
Empirical results show in many different way.
The evidence is least on the trade interaction which is the mechanism the theory relies on.
This paper examines the question for a group of countries that retained their currency in certain period.
Two research questions follow:

\textbf{Does exchange rate volatility against euro influence economic growth in the six non-euro European member states?}

\textbf{Does the effect of volatility on growth depend on the extent of country's trade integration with euro area?}